\section{Introduction}
\label{sec:intro}

The graph model has been widely used to denote relationships between different entities and to support data mining tasks across domains \cite{jiang2024spade+, ji2022detecting, li2024sefraud, liu2022user,niu2024discovering}. Application needs often necessitate \emph{inducing} composite relationships from heterogeneous data using meta-paths \cite{sun2011pathsim, fang2020effective, zheng2017entity} and representing them via \emph{relational graphs}. For example, on e-commerce platforms, market segmentation relies on community analysis in such \emph{induced} relational graphs that link customers with shared shopping preferences via meta-paths such as ``(\emph{customer}) $\rightarrow$ (\emph{product}) $\leftarrow$ (\emph{customer})''. Nevertheless, building such induced graphs is resource-intensive due to the cost of ``\textsf{join}'' operations on heterogeneous data~\cite{xirogiannopoulos2017extracting}. Moreover, although several studies have focused on accelerating relational graph materialization \cite{chatzopoulos2023atrapos, guo2024efficient, shi2016constrained, sun2011pathsim}, these methods are unsuitable for online business applications that require real-time on-demand analysis over time spans, rendering full materialization infeasible.

To address these challenges, we propose the analysis of induced relational graphs without full materialization. As a pioneer effort in this direction, we focus on \emph{hub queries} (\ced), which find important nodes in a relational graph according to a given \emph{centrality measure}.

\myp{Applications} \ced queries in relational graphs find applications in a variety of real-world scenarios, including:
\begin{itemize}[leftmargin=*]
    \item \etitle{Fraud Detection:} In an e-commence platform, by identifying hubs in relational graphs induced by ``(\emph{account}) $\rightarrow$ (\emph{transaction}) $\leftarrow$ (\emph{account})'' from transactional data, we can uncover financial accounts with an abnormally large number of interactions with other accounts, which might indicate the existence of money laundering \cite{cao2017hitfraud}.
    \item \etitle{Influence Analysis:} Hubs in a co-author network induced by ``(\emph{author}) $\rightarrow$ (\emph{paper}) $\leftarrow$ (\emph{author})'' \cite{mihara2015influence} denote influential researchers who are active in general domains. Moreover, hubs in a venue-filtered co-author network induced by meta-paths, e.g., ``(\emph{author}) $\rightarrow$ (\emph{paper}, \emph{venue}=`DB') $\leftarrow$ (\emph{author})'', reveal influencers in specific research areas \cite{shi2016constrained, li2014hrank}.
    \item \etitle{Protein Complex Discovery:} The protein-protein interaction (PPI) networks induced by meta-paths such as ``(\emph{protein}) $\rightarrow$ (\emph{pathway}) $\leftarrow$ (\emph{protein})'' from protein databases such as \textsf{STRING} (see \url{https://string-db.org/}) and \textsf{Reactome} (see \url{https://reactome.org/}) capture collective interactions among proteins in physiological pathways. Hubs in a PPI network are vital for the discovery of protein complexes \cite{wu2009core, srihari2013survey}.
\end{itemize}

\myp{Our Contributions}
We first consider \ced queries based on degree centrality. Inspired by cardinality sketches \cite{estan2003bitmap, alon1996space, durand2003loglog, heule2013hyperloglog, gibbons2001distinct}, we propose a novel sketch propagation framework for \ced queries on relational graphs with provable approximation guarantees. Specifically, we construct a \emph{neighborhood cardinality} sketch for each node to estimate its degree in the relational graph by iteratively propagating sketches along edges in matching instances of the given meta-path~$\mset$. Therefore, instead of computing the degree of each node after materialization, we collectively estimate the degrees of all nodes through a \emph{single} propagation process, which substantially improves the efficiency of processing \ced queries under degree centrality.

Then, we extend the sketch framework to approximate \ced queries based on the h-index \cite{LuZZS16}, which requires higher-order neighborhood information. Specifically, a key challenge is that the h-index of a node depends, in addition to its own degree, on the degrees of its neighbors, while cardinality sketches can only provide information on a node's own degree. To address this issue, we propose a \emph{pivot-based} algorithm built on the sketch propagation process. Instead of directly estimating h-indexes, we recursively find the nodes whose h-indexes reach or exceed that of a randomly chosen pivot node using sketches, in a quicksort-like fashion. As such, we can skip the nodes that cannot be hubs to process h-index-based \ced queries more efficiently. Note that we consider the $h$-index as a representative of centrality measures based on higher-order neighborhoods. By employing similar techniques, our framework can also be adapted to support other higher-order centrality measures, such as the i10-index and g-index \cite{egghe2006theory}. Furthermore, our framework can be extended to support closeness centrality \cite{olsen2014efficient, bergamini2019computing} by propagating cardinality sketches over the matching instances of $\mset$ in multiple iterations to estimate the number of nodes at different distances from each node in the relational graph.

Furthermore, we devise efficient pruning methods to answer \emph{personalized} \ced queries, that is, to determine whether a given node is a hub. Specifically, we maintain a lower bound on the number of nodes with centrality scores higher than the query node and terminate the sketch propagation process when that lower bound exceeds a given threshold. This early termination mechanism can accelerate personalized \ced queries based on both degree and h-index centrality measures.

Under the notion of $\epsilon$-separable sets, we show that the sketch propagation framework provides \emph{unbiased and efficient} approximations for \ced queries based on both centrality measures. The experiments reveal that the estimations converge much more quickly than the worst-case bound on the number of propagation iterations, achieving orders of magnitude speedups over exact methods.

Our main contributions are summarized as follows.
\begin{itemize}[leftmargin=*]
  \item We introduce the novel problem of hub queries (\ced) based on the degree and h-index centrality measures over relational graphs induced by meta-paths (Sect.~\ref{sec:preli}).
  \item We propose a \emph{sketch propagation} framework for approximate degree-based \ced queries and extend it to personalized \ced queries with early termination (Sect.~\ref{sec:scopeprop}).
  \item We devise a \emph{pivot-based} algorithm, also with early termination, for approximate (personalized) \ced queries using the h-index measure (Sect.~\ref{sec:hindex}).
  \item We show that our methods can scale well to massive graphs in which state-of-the-art exact methods fail to terminate in a reasonable time and mostly achieve orders of magnitude speedups over them while yielding accuracy beyond~90\%. We conduct case studies to indicate that our methods serve as seeding strategies comparable to state-of-the-art methods for influence analysis on relational graphs (Sect.~\ref{sec:exp}).
\end{itemize}
